\section{Asian and scope of the instrument}

The data set and evaluation protocol committed by the Intent Recognition Engine
(IRE, Component~1) are reproduced from the component report so that every figure
reported in \cref{ch:results} can be audited against the same artefact.

\section{Committed data set (1,500 records)}

The data set contains 1,500 evaluation records split into a well-formed set
(\(D^{+}\)) and a deliberately flawed set (\(D^{-}\)).

\begin{center}
\begin{tabular}{@{}lcc@{}}
  \toprule
  \textbf{Category} & \textbf{Records} & \textbf{Share}\\
  \midrule
  Well-formed (\(D^{+}\)) & 600 & 40.0\\
  Deliberately flawed (\(D^{-}\)) & 900 & 60.0\\
  \midrule
  Total \(D\) & 1,500 & 100.0\\
  \bottomrule
\end{tabular}
\end{center}

The eight flaw classes in the flawed subset are: adversarial utterances, RBAC
(role) violations, fabricated parameters, type mismatches, fabricated CVEs,
contradictory requests, out-of-scope targets, and ambiguous requests. Each
record carries a gold intent label and a gold flaw label, so per-record
decisions are scored deterministically.

\section{Measurement definitions}

The IRE defines three primary measures. Decision accuracy is the fraction of
records whose accepted/rejected decision matches the gold label (computed over
the well-formed set as acceptance accuracy and over the flawed set as catch
rate). The flaw-catch rate is the fraction of flawed records that are rejected.
The false-positive rate is the fraction of well-formed records that are wrongly
rejected. Decision accuracy is reported over the full 1,500-record set by
combining the acceptance rate on \(D^{+}\) and the catch rate on \(D^{-}\).

\section{Experiment schedule}

\begin{enumerate}
  \item \textbf{Experiment 1 (RQ1):} value of the guarded pipeline; raw model
  output against the full pipeline on the 1,500-record set; measures decision
  accuracy, catch rate, false-positive rate.
  \item \textbf{Experiment 2 (RQ2):} backend comparison; two offline backends
  (Gemma and Qwen) under an identical guardrail stack.
  \item \textbf{Experiment 3 (RQ2):} per-layer ablation; one protection removed
  at a time from the full pipeline, reporting the change in catch rate and
  false-positive rate relative to the baseline.
  \item \textbf{Experiment 4 (RQ2):} per-class catch rate; a deep-dive subset of
  the data set reporting the catch rate for each of the eight flaw classes.
\end{enumerate}

All results are single-run point estimates obtained on one CPU-only machine. The
evaluation harness commits per-record decisions so that the reported figures
can be reproduced.

\section{Reproducibility artefacts}

\begin{description}
  \item[Evaluation data set] -- the 1,500-record, gold-labelled collection used
  for all reported experiments.
  \item[Evaluation harness] -- the script that executes the four experiments and
  writes per-record decisions.
  \item[Per-layer ablation protocol] -- the configuration used to remove one
  protection at a time.
  \item[Test suite] -- 506 tests (434 unit, 72 integration) collected on
  7 October 2026, all passing, with 76\% statement coverage.
\end{description}